\title{Large Language Models Can Automatically Engineer Features for Few-Shot Tabular Learning}

\begin{abstract}
Large Language Models (LLMs) have demonstrated exceptional proficiency in solving complex and novel reasoning tasks, revealing significant promise for applications in tabular learning that are essential across numerous practical fields. In this work, we introduce a new in-context learning approach, \textsf{FeatLLM}, which leverages LLMs to act as feature engineers, constructing an input dataset that is optimally configured for predictive modeling over tabular data. These newly created features are subsequently utilized by a straightforward machine learning model—such as linear regression—for the inference of class probabilities, enabling effective few-shot learning performance. At inference, our \textsf{FeatLLM} paradigm relies exclusively on this simple prediction model in conjunction with the engineered features. Unlike previously established LLM-based strategies, \textsf{FeatLLM} does not require the LLM to be invoked at prediction time for individual samples. Additionally, it only necessitates API-level interaction with the LLM and circumvents any issues related to prompt length restrictions. Through comprehensive experimentation on a variety of tabular datasets from distinct application domains, we demonstrate that \textsf{FeatLLM} constructs high-quality, rule-based features, surpassing existing competitors such as TabLLM and STUNT by an average margin of 10\%.
\end{abstract}

\section{Introduction}
Recent advances have illustrated that Large Language Models (LLMs) excel at generating appropriate outputs for tasks they have not previously encountered, all without the need for fine-tuning tailored to individual tasks~\cite{creswell2022selection,imani2023mathprompter,taylor2022galactica}. Trained on vast collections of textual data, these models encapsulate a wealth of prior knowledge, which can often serve in place of specialized domain expertise across multiple fields~\cite{Genomes2021,singhal2023large,wilcox2003role}. Their utility spans a spectrum of automated applications, including but not limited to, dialogue understanding~\cite{finch2023leveraging} and software bug fixing~\cite{fan2023automated}. Notably, by presenting LLMs with explicit task descriptions and a handful of illustrative examples in their input prompts, these models are able to grasp contextual cues for a given task and identify salient features or instances~\cite{brown2020language,zhao2023survey}. This capability suggests that LLMs can function effectively in analyzing structured data, indicating their suitability for sophisticated feature engineering tasks.

Extending the use of LLMs’ internalized knowledge and logical reasoning to tabular data analysis has gained attention in recent studies. For example, leveraging knowledge about demographic risks—such as the predisposition of older adults to certain illnesses—can greatly enhance predictive modeling in healthcare. To harness this potential, prior research has transformed tasks and feature information into natural language, reformatted raw tables, and provided this information to LLMs~\cite{dinh2022lift,wang2023anypredict}. These frameworks then commonly employ in-context learning~\cite{nam2023semi} or rely on efficient fine-tuning mechanisms~\cite{dinh2022lift,hegselmann2023tabllm}. Empirical evidence supports that integrating LLM-driven prior knowledge significantly boosts performance, especially in settings with limited training examples, consistently outperforming conventional approaches to tabular learning in these scenarios~\cite{hegselmann2023tabllm}.

Existing LLM-based methods for tabular data learning encounter several significant limitations. When performing end-to-end predictions, each individual sample necessitates at least one full LLM inference, resulting in high computational overhead. Moreover, achieving optimal predictive performance usually depends on fine-tuning the LLM, a process that is impractical if the base LLM lacks accessible training code or public tuning APIs. This issue is especially pertinent given that many state-of-the-art LLMs, as noted in recent work~\cite{anil2023palm,openai2023gpt4}, only permit restricted interactions via inference APIs. Additionally, most available strategies are ill-equipped to handle extended input prompts~\cite{liu2023lost}; as the feature count in tabular datasets increases and the serialized input text grows longer, these approaches either become unusable or show noticeable performance degradation. Collectively, such barriers significantly impede the practical adoption of LLM-based tabular learning techniques.

To address these challenges, our method shifts from utilizing LLMs for direct, end-to-end predictions toward employing them as tools for automated feature construction. We present a new in-context learning paradigm, \textsf{FeatLLM}, which capitalizes on the LLM’s intrinsic ability to distill the “criteria” underlying its predictions. For example, given a disease classification task, the LLM can be prompted to articulate explicit decision rules based on feature configurations that characterize the disease. These derived criteria or rules are then leveraged to generate novel features, effectively substituting for the original set in downstream tasks. This realignment harnesses the LLM’s strengths for feature derivation, thereby enabling the deployment of a straightforward downstream model; after the relevant features are extracted, there is no longer a necessity for complex models during inference, resulting in faster response times. By exploiting in-context learning and supplying illustrative examples within prompts, feature extraction is accomplished without the burdens of model retraining or fine-tuning. To further address issues related to excessive prompt length, we incorporate ensemble learning strategies~\cite{breiman2001random}, such as feature bagging, to manage input size effectively.

\textsf{FeatLLM} organizes the task of prompt-based feature engineering into two main components: (1) interpreting the problem statement and discerning the linkage between input features and target variables; and (2) utilizing insights obtained from the initial analysis, alongside limited annotated examples, to formulate discriminative rules capable of separating different classes. This systematic reasoning directs LLMs to ignore features that do not contribute useful information for rule generation. The resulting rules are subsequently implemented on the dataset (with parsing as code), converting the original data into binary features that indicate compliance with each rule on a per-sample basis. A simple, low-capacity predictive model is then trained to map these binary features to class probabilities. To further enhance reliability, \textsf{FeatLLM} leverages an ensemble mechanism, repeatedly constructing features using a feature bagging procedure. By carefully controlling the number of selected features and samples per bag, the method addresses the practical challenge of overly large input prompts. As \textsf{FeatLLM} requires no additional model training and incurs minimal inference overhead, large language models can be efficiently leveraged for a broad array of real-world applications.

The effectiveness of \textsf{FeatLLM} is demonstrated on a suite of 13 tabular datasets within low-shot learning contexts, where it achieves notable accuracy and stability. Empirical results show that LLMs can effectively draw on both their inherent knowledge and patterns present in the data itself, leading to the creation of informative and predictive features. Our proposed system consistently surpasses other state-of-the-art few-shot learning methods across multiple benchmarks. Implementation resources are publicly accessible via an anonymized GitHub repository at~\texttt{https://github.com/Sungwon-Han/FeatLLM}.

\section{Related Work}

\subsection{Few-Shot Learning with Tabular Data} The progression of few-shot learning techniques has made it feasible to acquire broadly applicable representations from datasets while incurring minimal annotation expense~\cite{chen2018closer}. Although early research in this domain centered around image data, the applicability of few-shot learning has successfully expanded to additional modalities such as text, audio, and more recently, tabular formats~\cite{majumder2022few,nam2023stunt,schick2021exploiting}. A primary challenge when working with limited labeled examples lies in the susceptibility to erroneous correlations~\cite{bartlett2021deep}. To counteract this, more recent methodologies have integrated auxiliary data resources or leveraged domain-relevant inductive biases by incorporating prior knowledge during model development. One prominent strategy involves leveraging large-scale unlabeled tabular datasets for semi-supervised learning enhanced by targeted data augmentation~\cite{ucar2021subtab,yoon2020vime}, or establishing robust model parameters through broad, unsupervised pre-training not specific to a particular task~\cite{somepalli2021saint}. These unsupervised paradigms often utilize data manipulation such as masking or corrupting input portions, followed by a reconstruction loss~\cite{arik2021tabnet,majmundar2022met}, or employ data augmentation in conjunction with contrastive objectives~\cite{bahri2022scarf,chen2020simple}. Nevertheless, due to the inherently diverse structures in tabular data, designing augmentation techniques that generalize well is nontrivial, resulting in these approaches demonstrating limited success in few-shot contexts~\cite{nam2023stunt}. 

A line of recent studies investigates transfer learning by pre-training models on an assortment of tabular datasets unrelated to the target, and then adapting to the specific target task~\cite{zhang2023generative,levin2022transfer,zhu2023xtab}. One example is TransTab~\cite{wang2022transtab}, where transformer-based architectures are trained to learn inter-column semantic relations utilizing both column names and table entries. This form of representation transfer has been shown to improve the performance for previously unseen tabular tasks, enabling zero-shot or few-shot applications. Another approach, exemplified by TabPFN~\cite{hollmann2022tabpfn}, pre-trains a Bayesian neural network using synthetic datasets generated from mixtures of data distributions. During evaluation, the model operates by processing the complete set of training and test instances collectively, without further fine-tuning. Knowledge attained from this wide-ranging pre-training led to outcomes surpassing traditional tree-based techniques, demonstrating particular effectiveness in low-shot learning scenarios.

\subsection{Language-Interfaced Tabular Learning} The use of large language models (LLMs)~\cite{anil2023palm,openai2023gpt4} represents a further direction for injecting prior knowledge into tabular learning frameworks. Having been trained on vast, heterogeneous collections of textual data, LLMs have shown impressive transfer capabilities, adapting to novel tasks within a variety of fields~\cite{brown2020language}. Recent research efforts have explored how LLMs' embedded prior knowledge can benefit tabular learning by representing table data as serialized text for inclusion in model prompts~\cite{dinh2022lift,hegselmann2023tabllm,wang2023anypredict}. These approaches encode the table structure, feature metadata, and task context into the input, facilitating the inference of feature-task relationships by the LLM. There is an emerging trend toward either parameter-efficient training of language models—via mechanisms such as LoRA~\cite{hu2021lora} or IA3~\cite{liu2022few}—or towards harnessing in-context learning by introducing few-shot exemplars in prompts, obviating the need for further parameter updates~\cite{dinh2022lift,hegselmann2023tabllm,nam2023semi,slack2023tablet}.

Although the aforementioned approaches have demonstrated strong performance in enhancing few-shot tabular learning, their dependency on routing all inference through an LLM creates practical obstacles for deployment in industry settings. This paper seeks to break from the standard end-to-end black-box paradigm in which every prediction flows through an LLM. Rather than evaluating individual instances through the LLM, our approach directs the LLM to explicitly deduce the governing conditions or rules that define each class label. \textsf{FeatLLM} then interprets these LLM-generated rules as programmatic code, which is subsequently used to engineer new features. This enables inference that bypasses the LLM, taking advantage of in-context learning to extract rules informed both by prior knowledge and a limited set of examples.

\section{Methods}
The primary objective of \textsf{FeatLLM} is to elicit a set of ``rules"—that is, attribute-based conditions linked to each potential class label—from the provided few-shot samples, as illustrated in Figure~\ref{fig:main_model}. Once these rules are formulated with the aid of an LLM, they are transformed into new binary features for each data point, each denoting compliance with a particular rule. These resultant binary feature vectors are then aggregated using a dot product with non-negative, trainable weights in order to quantify the probability that a given sample belongs to each class. During model training, the procedure allows the relative influence of each engineered feature on class membership to be learned (Section~\ref{Sec:training}). This pipeline is reiterated across multiple bootstrapped sets to enable bagging, with the outcomes integrated through ensembling. The specifics of the problem setup as well as a breakdown of the individual stages are presented in the following subsections.

\vspace*{-0.12in}
\paragraph{Problem formulation.} Suppose we have a tabular dataset comprising $N$ labeled instances, represented as $\mathcal{D} = \{(\mathbf{x}^i, \mathbf{y}^i)\}_{i=1}^{N}$. Each sample $\mathbf{x}^i$ is a $d$-dimensional input vector, reflecting $d$ distinct input features. The dataset $\mathcal{D}$ is accompanied by natural language identifiers for these features, written as $F = \{f_j\}_{j=1}^{d}$, with corresponding standalone definitions. In a classification context, $\mathbf{y}^i$ signifies membership in a set of predetermined classes. For $k$-shot learning scenarios, the model is trained only on a random selection of $k$ ($<N$) labeled examples from the dataset.

\subsection{Prompt Design for Extracting Rules}
\label{Sec:prompt_design}
To enable LLMs to derive rules using more precise logical steps, we structure the problem-solving method to reflect the typical reasoning process of a human expert dealing with tabular data. The designed prompt incorporates three primary elements (illustrated in Fig.~\ref{fig:prompt_for_rule} and detailed in Appendix~\ref{sec:example_rule}):

\vspace*{-0.12in}
\paragraph{Basic information description.} This section conveys fundamental details necessary for addressing the task at hand (refer to the orange-highlighted portion in Fig.~\ref{fig:prompt_for_rule}). It encompasses a concise explanation of the task, specifying label details, together with feature summaries and practical examples. The task introduction is articulated as an inquiry—for instance, ``Does this patient's myocardial infarction complications data indicate chronic heart failure? Yes or no?" (see additional instances in Appendix~\ref{sec:task_desc}). Feature summaries identify the type of value and, where applicable, include the feature's definition. When dealing with categorical features, samples of potential categories may also be specified. For illustrative purposes, select training cases are linearized as textual inputs alongside their correct labels, using the serialization template established in prior work~\cite{dinh2022lift,hegselmann2023tabllm}:

\begin{align}
&\text{Serialize}(\mathbf{x}^i,\mathbf{y}^i, F) = \nonumber \\ 
&``\ f_1\ \text{is}\ \mathbf{x}^i_1.\ ... \ f_d\  \text{is}\  \mathbf{x}^i_d.\ \ \text{Answer:}\  \mathbf{y}^i\ ”
\end{align}

\vspace*{-0.12in}
\paragraph{Reasoning instruction.} Our approach to improving the LLM's reasoning capabilities involves the use of tailored reasoning instructions (see blue text in Fig.~\ref{fig:prompt_for_rule}). Mirroring the structure of the chain-of-thought methodology~\cite{wei2022chain}, these instructions start with a brief introductory statement. Subsequently, we segment the reasoning procedure of the LLM into two distinct stages. In the first stage, the LLM is prompted to analyze and infer potential causal relationships or underlying patterns connecting the features and the task description. This reasoning occurs in the absence of illustrative examples, encouraging the LLM to draw upon its general background knowledge and intuitive understanding, and thereby allowing it to tap into its latent knowledge base concerning the problem context. In the subsequent stage, the LLM incorporates example demonstrations, building upon the insights gained in the previous step, to formulate rules pertinent to each class. This bifurcated process mitigates the risk of the model learning unintended correlations from irrelevant attributes and encourages greater attention toward the more informative features. To avoid the risk of extracting an insufficient or an excessive number of rules—which could lead to underfitting or result in overly lengthy prompts—a uniform rule count (i.e., 10) is designated for extraction per class.\footnote{Refer to Section~\ref{sec:analysis} for detailed discussion about the rule count selection.} Furthermore, to ensure that the rules correspond appropriately to their respective feature types, the LLM is instructed to utilize the basic information section's feature type listings when constructing rules, thereby reducing the likelihood of type inconsistency.

\vspace*{-0.12in}
\paragraph{Response instruction.} In order to streamline the process of parsing and leveraging the rules produced by the LLM, we incorporate detailed response formatting guidelines (see yellow text in Fig.~\ref{fig:prompt_for_rule}). The prompting includes explicit specifications regarding the anticipated structure for responses corresponding to each answer class (i.e., required response format) and prescribes the layout that each rule should adhere to (i.e., required rule format). This direction equips the LLM to select an appropriate template aligned with the nature of the feature involved. Where further instruction is not imposed, the LLM is permitted to utilize logical connectives, such as AND and OR, when composing combined rules. Representative outputs are detailed in Appendix~\ref{sec:outcome-rule}.

\subsection{Inferring Class Likelihood via Rules}
\label{Sec:training}

\paragraph{Parsing rules for feature generation.} The rules produced in the prior phase are utilized to construct new binary features. For each class, these features encode whether the instance adheres to its corresponding rule set, taking values of either '0' or '1.' Such features provide a basis for estimating the probability that a given sample belongs to a specific class. However, as the LLM produces rules in natural language, a parsing mechanism must be established to translate these rules automatically into binary features. Although rule response formats are constrained to facilitate parsing, LLM outputs can still exhibit syntactic inconsistencies~\cite{kovavc2023large}. For example, the model may replace operators such as ``$>$” or ``$<$” with verbal counterparts like ``greater than" or ``less than," thereby complicating automatic parsing.

To overcome the issue of variable, potentially noisy text, we opt against designing intricate parsers; instead, we repurpose the LLM for parsing. LLMs are adept at interpreting meaning regardless of textual variations, and they possess the capability to write concise code from specifications, owing to their training on large code corpora. To exploit these strengths, we formulate prompts containing function signatures, along with input and output specifications, and the derived rules, submitting these to the LLM. Illustrative prompts and corresponding outputs can be found in Figures~\ref{fig:prompt_for_parsing} and~\ref{sec:example_parsing}-\ref{sec:example_parse_outcome} in the Appendix. The resulting code fragments are then run via Python's \texttt{exec()} to process and convert the data.

\vspace*{-0.12in}
\paragraph{Inferring class likelihood.} Once these rule-based binary features are generated per class, an elementary approach to estimating the likelihood that a sample belongs to a particular class is to tally the number of rules satisfied for each class, i.e., summing the respective binary features. Nevertheless, features and rules may differ in their predictive value, requiring that their relative importance be inferred from data. For this, we utilize standard machine learning (ML) algorithms, assigning a model to each set of class-wise binary features. One suitable choice is a bias-free linear model. Formally, let $\mathbf{z}_k^i \in \{0, 1\}^R$ denote the binary vector for sample $\mathbf{x}^i$ and class $k$, where $R$ represents the total rules per class and $c$ is the total number of classes. The class probability vector $\mathbf{p}^i$ can then be obtained as:

\begin{align}
\text{logit}_k^i &= \max(\mathbf{w}_k, 0) \cdot \mathbf{z}_k^i, \nonumber \\
\mathbf{p}^i &= \text{Softmax}({[\text{logit}_{1}^i, ..., \text{logit}_{c}^i]}). \label{eq:infer_prob}
\end{align}

In this framework, $\mathbf{w}_k$ denotes the set of learnable weights within the linear model, indicating the relative importance attributed to each feature. Restricting these weights to a positive domain guarantees that features produced by the LLM cannot decrease a class's probability during the learning stage. Afterward, the resulting logit vector undergoes the softmax transformation to yield class probabilities. The weights $\mathbf{w}_k$ are optimized by minimizing the cross-entropy loss with a limited selection of training examples.

\vspace*{-0.12in}
\paragraph{Ensembling with bagging.} The entire workflow is iteratively executed to construct multiple inference models, ultimately combining them via ensembling to reach a final decision. Given $T$ trained weights, $\{\mathbf{w}[t]\}_{t=1}^T$, predictions are determined by calculating the mean of per-class probabilities, where each $\mathbf{w}[t]$ produces a probability vector $\mathbf{p}^i$, as formalized in Eq.~\ref{eq:infer_prob}.

To ensure sufficient variation among generated rules for ensemble modeling, LLM inference is performed using a relatively high temperature. Additionally, the order of few-shot demonstration examples is randomized, given prior findings that permutation can influence the LLM's responses~\cite{lu2022fantastically}\footnote{Further discussion of rule diversity appears in Appendix~\ref{sec:diversity}}. Exploiting the ensemble's diversity to achieve robust results, bagging is employed to randomly select a subset of features or samples for each iteration, which becomes especially beneficial as the quantity of data or features grows. The ensembling strategy brings two primary benefits. First, if certain LLM-generated rules are erroneous in some trials, their influence can be offset by correct rules in other iterations, thus increasing the overall robustness of the model. This is related to the self-consistency property of LLMs~\cite{wang2022self}, since rules that consistently emerge across repeated runs are more likely to be reliable. Second, ensembling alleviates the LLM prompt length constraint. For cases where the tabular input surpasses permissible prompt dimensions, bagging effectively reduces input size while safeguarding accuracy. Although a single rule set may overlook relationships involving some features, aggregating diverse weak learners from multiple trials compensates for any omissions, collectively ensuring comprehensive feature and instance representation.

\section{Experiments}
We assess the effectiveness of \textsf{FeatLLM} across a diverse selection of tabular datasets and carry out comprehensive studies to understand the underlying mechanisms of our proposed approach. Owing to limited space, further experiments—including those that explore different LLM variants, conduct in-depth qualitative analyses, and additional investigations—are provided in the Appendix.

\subsection{Performance Evaluation}
\paragraph{Datasets.}
Our evaluation covers 13 different datasets spanning binary and multi-class classification tasks: (1) Adult~\cite{asuncion2007uci}, used to determine if an individual's yearly income surpasses \$50,000; (2) Bank~\cite{moro2014data}, focused on predicting term deposit subscriptions among customers; (3) Blood~\cite{yeh2009knowledge}, aimed at forecasting whether blood donors will donate again; (4) Car~\cite{kadra2021well}, concerning vehicle quality assessment; (5) Communities~\cite{misc_communities_and_crime_183}, for predicting crime rates in particular communities; (6) Credit-g~\cite{kadra2021well}, which classifies individuals as having either high or low credit risk; (7) Diabetes\footnote{\texttt{kaggle.com/datasets/uciml/pima-indians-diabetes-database}}, tasked with diagnosing diabetes in individuals; (8) Heart\footnote{\texttt{kaggle.com/datasets/fedesoriano/heart-failure-prediction}}, for coronary artery disease prediction; and (9) Myocardial~\cite{misc_myocardial_infarction_complications_579}, focused on identifying cases of chronic heart failure.

Additionally, our study includes datasets that are either newly released or artificially generated—these have received less attention in prior literature and were not accessible during the pre-training phases of LLMs: (10) Cultivars, which estimates the potential grain yield for a specific soybean cultivar\footnote{\texttt{archive.ics.uci.edu/dataset/913}}; (11) NHANES, which predicts a person's age group based on the recorded attributes\footnote{\texttt{archive.ics.uci.edu/dataset/887}}; (12) Sequence-type, assigning each input sequence to one of four categories: arithmetic, geometric, Fibonacci, or Collatz; and (13) Solution-mix, which decides if a mixture's percent concentration—derived from four constituent solutions—exceeds 0.5. The first two recent datasets were added to the UCI Machine Learning Repository after September 2023, guaranteeing their absence from the gpt-3.5-turbo-0613 LLM’s pre-training corpus. The last two datasets are synthetically constructed, ensuring that they remain unfamiliar to the LLM prior to evaluation.

These datasets present a wide range of sizes and structural intricacies, which are itemized in Table~\ref{Tab:basic-desc}. Attribute names and dataset descriptions are unambiguous and thoroughly defined. Datasets such as `Communities' and `Myocardial' comprise over 100 attributes, challenging the LLM's ability to accommodate such breadth within a restricted context window.

\vspace*{-0.12in}
\paragraph{Baselines.}
For benchmarking, we compare \textsf{FeatLLM} against nine baseline models. The initial three consist of established tabular data classification techniques: (1) Logistic Regression (LogReg), (2) XGBoost~\cite{chen2016xgboost}, and (3) RandomForest~\cite{ho1995random}. Two subsequent baselines—(4) SCARF~\cite{bahri2022scarf} and (5) STUNT~\cite{nam2023stunt}—operate under the assumption that copious amounts of unlabeled data are available; it is worth noting, however, that this condition might not always hold in practice. Another contemporary model, (6) TabPFN~\cite{hollmann2022tabpfn}, leverages synthetic datasets generated from varied distributions for model pretraining. The last set comprises methods reliant on LLMs: (7) In-context learning (In-context)~\cite{wei2022emergent}, (8) TABLET~\cite{slack2023tablet}, and (9) TabLLM~\cite{hegselmann2023tabllm}. The In-context learning approach injects a small number of labeled examples directly into the prompt for inference without any fine-tuning. TABLET enhances inference by integrating supplemental prompt hints using rule sets and prototype-based guidance derived from an external classifier. TabLLM adopts parameter-efficient finetuning, such as IA3~\cite{liu2022few}, allowing LLMs to

\vspace*{-0.12in}
\paragraph{Implementation details.}
\textsf{FeatLLM} employs GPT-3.5 as the foundational large language model, though its architecture is compatible with alternative LLM backbones (refer to Appendix~\ref{sec:backbones} for experiments using various backbones). During LLM inference, a temperature of 0.5 is adopted, and the API top-p parameter is maintained at its default of 1. For ensemble creation and rule extraction, the model utilizes 20 ensembles and identifies 10 rules, respectively. Additional analysis on hyper-parameter selection can be found in Figure~\ref{fig:hyperparameter2} and Appendix~\ref{sec:hyper-parameter}. The linear model is trained using the Adam optimizer with a learning rate set at 0.01, running for 200 epochs. For selecting the best epoch, $k$-fold cross-validation is applied. When re-implementing comparison methods, in-context learning is conducted using GPT-3.5, while for TabLLM the T0~\cite{sanh2022multitask} model is used to match its original configuration. It is important to note that TabLLM limits permissible LLMs to those with publicly available checkpoints, thus GPT-3.5 cannot be employed in those scenarios. Conventional machine learning approaches, including LogReg, XGBoost, and RandomForest, have their parameters fine-tuned via Grid-search in conjunction with $k$-fold cross-validation. The parameter $k$ is chosen as either 2 or 4 to ensure that every class is represented in the training data for each split. More comprehensive implementation specifics are available in Appendix~\ref{sec:baselines}.

\vspace*{-0.12in}
\paragraph{Results.}
Tables~\ref{tab:main1} and \ref{tab:main2} show comparative results spanning multiple datasets. For each setting, we conduct experiments three times, altering the random division of training and test datasets, and report the mean as well as standard deviation of AUC scores. Across benchmarks, \textsf{FeatLLM} either leads or places second relative to other baselines. Remarkably, our framework, using only 4 shots, attains performance on par with classic tabular models that utilize 32 shots (as illustrated in Fig.~\ref{fig:compares}). While STUNT and TabPFN displayed notable gains on select benchmarks, their aggregate improvements over standard machine learning approaches were modest. In addition, approaches leveraging LLMs, such as in-context learning, TABLET, and TabLLM, proved incapable of handling tabular data containing a large feature set. In contrast, our method, leveraging both bagging and ensembling strategies, is robust even on datasets with high-dimensional features and achieves strong results. It should be emphasized that our bagging and ensemble methodology can potentially be integrated with other LLM-based techniques. However, doing so would effectively require double the per-sample inference workload, significantly increasing computational demands and rendering such approaches infeasible in practice.

\subsection{Analysis \& Discussion}
\label{sec:analysis}
\paragraph{Ablation study.} To dissect the influence of critical components, we conduct a suite of ablation experiments targeting the following: (1) \textsf{-Tuning}: eliminates the linear model’s weight tuning step, instead aggregating the binary feature to derive the logit; (2) \textsf{-Ensemble}: omits ensemble aggregation; (3) \textsf{-Description}: removes inclusion of feature descriptions; and (4) \textsf{-Reasoning}: leaves out the initial reasoning instruction (Step-1) during rule generation.

The results in Table~\ref{tab:ablation} report average AUC changes across datasets due to these ablations. Each removal invariably leads to diminished performance. Notably, the utility of weight tuning increases with greater numbers of shots, suggesting that with more data, the system can better discern rule significance and thus benefit more from tuning. Conversely, both feature descriptions and reasoning guidance contribute most when limited data is available, highlighting the importance of effectively leveraging LLM’s prior knowledge under few-shot conditions. Finally, the ensemble approach proposed here delivers consistent improvements across all settings.

\vspace*{-0.12in}
\paragraph{Training \& inference time comparison.} We evaluate the training and inference durations across various methods. Our assessment utilizes the Adult dataset with a training set comprising 16 samples. All methods, with the exception of TabLLM (which leverages four A100 GPUs for model parallelization), are executed on a single A100 GPU by default. Table~\ref{tab:cost} presents a summary of these time comparisons. Importantly, \textsf{FeatLLM} achieves inference times on par with traditional baseline methods. In sharp contrast, approaches reliant on LLMs—namely In-context learning, TABLET, and TabLLM—suffer from significantly increased inference times. When considering training time, \textsf{FeatLLM} matches other LLM-based baselines: it requires only 30 API queries, imposing minimal cost and permitting efficient parallelization.

\vspace*{-0.12in}
\paragraph{Quality of features generated by \textsf{FeatLLM}.}
To assess the utility of the rule-based features produced by \textsf{FeatLLM} in practical applications, we run a straightforward experiment. This involves calculating the average AUROC between each generated feature and the target outcome, followed by determining the proportion of features attaining an AUROC greater than 0.5 across each dataset. The outcomes, reported in Table~\ref{tab:quality}, indicate that most generated rules are pertinent to the prediction target and furnish valuable information for the learning task (see the ``Before tuning'' column). Furthermore, during subsequent linear model tuning, rules exhibiting weak associations with the data—identified by learned weights beneath a specified threshold—are systematically excluded (refer to ``After tuning''). Here, the threshold is set to 0.1, reflecting the overall weight distribution.

\vspace*{-0.12in}
\paragraph{Effect of adding spurious correlations.} To evaluate how different approaches mitigate the influence of noisy spurious correlations in few-shot learning scenarios, we performed an investigation. Our methodology involved the Heart dataset, augmented by appending randomly selected columns from the Adult dataset—columns that bear no relevance to the primary predictive task of the Heart dataset. We tracked how each model's performance varied as more extraneous columns were introduced. For consistency across approaches, we restricted all methods to utilize just 8 labeled samples, omitted any use of unlabeled data, and thus did not include approaches such as SCARF or STUNT in our analysis.

Figure~\ref{fig:spurious} displays the resulting changes in model performance in the presence of additional spurious correlations. Notably, \textsf{FeatLLM} suffered the smallest drop in performance compared to other baselines. This robustness appears attributable to our framework’s reliance on prior knowledge that guides the identification of meaningful relationships between features and the predictive task when formulating rules, thereby largely avoiding spurious rules linked to unrelated columns. Empirically, rules involving noisy, unrelated columns made up merely 1-2\% of all extracted rules. Nevertheless, it is important to note that \textsf{FeatLLM} remains susceptible to incorporating spurious correlations and potentially misconstruing causality between features and the prediction target when random columns are derived from data sources closely related to the Heart dataset (refer to Appendix~\ref{sec:spurious-appendix} for more details).

\vspace*{-0.12in}
\paragraph{Hyper-parameter analysis}
\textsf{FeatLLM} relies primarily on two hyper-parameters: the ensemble size and the count of rules per class obtained per LLM invocation. To evaluate the influence of these hyper-parameters, we performed a series of experiments. Our findings reveal that increasing the ensemble size leads to better outcomes, but that improvement plateaus when the ensemble reaches approximately 20 models (refer to Fig.~\ref{fig:appendix-hyperparameter1} in the Appendix). Regarding the rule count, we did not observe statistically significant variations across settings; however, choosing 10 rules strikes an optimal balance between prompt length and prediction quality (refer to Fig.~\ref{fig:hyperparameter2}). We hypothesize that expanding the number of rules increases the dimensionality of the input, thus introducing additional information, but may also prompt overfitting, especially under low-shot conditions.

\section{Conclusion}
We introduce \textsf{FeatLLM}, an approach that pushes the boundaries of few-shot tabular learning with LLMs. Instead of employing LLMs exclusively for instance-level predictions, \textsf{FeatLLM} adopts a feature engineering perspective—utilizing LLMs to produce predictive rules. These rules are aggregated via an ensemble and bagging methodology, allowing \textsf{FeatLLM} to achieve robust predictive accuracy with minimal inference overhead, requiring only API-level LLM access and not additional training, while also addressing limitations related to input feature dimensions. Altogether, \textsf{FeatLLM} demonstrates high accuracy and presents an efficient alternative for leveraging LLMs on practical tabular problems.

Although \textsf{FeatLLM} is currently tailored for low-shot learning applications, our future efforts aim to extend its functionality for larger datasets by generating new kinds of features beyond rules, enhancing interpretability, and enabling broader applicability to various tasks.

\section{Broader Impact} The \textsf{FeatLLM} framework is designed to harness LLMs’ stored knowledge in order to enhance tabular prediction tasks, moving beyond the capability of standard supervised tabular learning systems. This work contributes to improved data efficiency by leveraging prior knowledge, which helps mitigate overfitting that often arises with scarce training data. \textsf{FeatLLM} is especially advantageous for real-world use cases—in sectors such as finance or healthcare—where annotations are expensive or hard to procure, and the extensive knowledge embedded within LLMs can be leveraged (given pretraining on finance- or healthcare-related texts). For future work aimed at broader deployment, it will be critical to investigate and understand the potential propagation of societal biases and misinformation embedded within LLM prior knowledge, since this prior may substantially influence, or even overshadow, predictions relative to available labeled examples.

\end{document}